\section*{TÓM TẮT}
\phantomsection\addcontentsline{toc}{section}{TÓM TẮT}
Những tiến bộ gần đây trong học sâu đã nâng cao đáng kể khả năng nén ảnh bằng cách cho phép các mô hình học tập vượt trội so với các tiêu chuẩn truyền thống trong cả hiệu suất tỷ lệ--độ méo và chất lượng nhận thức. Đặc biệt, các phương pháp dựa trên Mạng Đối Thủ Sinh Sinh (GAN) như Nén Ảnh Độ Trung Thực Cao (HiFiC) đã chứng minh khả năng mạnh mẽ trong việc bảo toàn các chi tiết hình ảnh tốt nhất ở tỷ lệ bit thấp. Tuy nhiên, độ phức tạp tính toán cao của các mô hình dựa trên GAN gây ra những thách thức lớn cho việc triển khai thời gian thực và các hệ thống nhúng.

Luận văn này trình bày một phương pháp thiết kế phần cứng--phần mềm đồng thời cho nén ảnh học tập dựa trên mô hình HiFiC GAN được lượng tử hóa. Hệ thống được đề xuất nhắm vào giai đoạn suy luận của đường dẫn HiFiC, trong đó các thành phần tính toán yêu cầu nhiều được tăng tốc trên FPGA, trong khi các chức năng còn lại được thực hiện trên phần mềm. Để cho phép triển khai phần cứng hiệu quả, mạng nơ-ron được lượng tử hóa thành độ chính xác FP16 và được triển khai bằng Tổng hợp cấp cao (HLS), cho phép khai thác hiệu quả tính song song và phát triển thiết kế hợp lý.

Kiến trúc bộ tăng tốc áp dụng mô hình dòng dữ liệu truyền phát với các yếu tố xử lý song song để ánh xạ hiệu quả các phép toán phức tạp vào phần cứng. Thiết kế được triển khai và đánh giá trên nền tảng FPGA ZCU104, chứng minh tần số hoạt động khả thi, sử dụng tài nguyên chấp nhận được và độ chính xác số tương đương với kết quả tham chiếu phần mềm.

Kết quả thực nghiệm cho thấy rằng khung thiết kế phần cứng--phần mềm đồng thời được đề xuất cung cấp một giải pháp hiệu quả và thực tiễn để tăng tốc các mô hình nén ảnh học tập. Công việc này nhấn mạnh tiềm năng của tăng tốc dựa trên FPGA cho nén ảnh dựa trên GAN nâng cao và đóng vai trò là nền tảng cho các mở rộng trong tương lai hướng tới triển khai phần cứng toàn diện hơn.
\cleardoublepage

\section*{ABSTRACT}
\phantomsection\addcontentsline{toc}{section}{ABSTRACT}
Recent advances in deep learning have significantly improved image compression by enabling learned models to outperform traditional hand-crafted standards in both rate--distortion efficiency and perceptual quality. In particular, Generative Adversarial Network (GAN)-based approaches such as High-Fidelity Image Compression (HiFiC) have demonstrated strong capability in preserving fine visual details at low bitrates. However, the high computational complexity of GAN-based models poses major challenges for real-time deployment and embedded systems.

This thesis presents a hardware--software co-design approach for learned image compression based on a quantized HiFiC GAN model. The proposed system targets the inference stage of the HiFiC pipeline, where computationally intensive components are accelerated on FPGA, while the remaining functions are executed in software. To enable efficient hardware implementation, the neural network is quantized to FP16 precision and implemented using High-Level Synthesis (HLS), allowing effective exploitation of parallelism and streamlined design development.

The accelerator architecture adopts a streaming dataflow model with parallel processing elements to efficiently map convolutional operations onto hardware. The design is implemented and evaluated on the ZCU104 FPGA platform, demonstrating feasible operating frequency, acceptable resource utilization, and numerical accuracy comparable to software reference results.

Experimental results indicate that the proposed hardware--software co-design framework provides an effective and practical solution for accelerating learned image compression models. This work highlights the potential of FPGA-based acceleration for advanced GAN-based image compression and serves as a foundation for future extensions toward more comprehensive hardware deployment.
\cleardoublepage